\section{Related Work}
\textbf{Video Diffusion Models}
Video diffusion models have advanced rapidly, beginning with UNet~\citep{ronneberger2015u} based approaches that extended image diffusion backbones into the temporal domain~\citep{ho2022imagen,ho2022video,he2022latent,blattmann2023stable}. These early models enabled short-form video generation but faced limitations in scalability. The introduction of the Diffusion Transformer (DiT)\citep{peebles2023scalable} represented a turning point by replacing convolutional hierarchies with transformer blocks, which allowed models to capture global spatio-temporal dependencies more effectively and to scale with larger datasets and computational resources. This shift led to a new wave of architectures, such as Sora\citep{openai2024sora}, which produces realistic, coherent videos with strong temporal consistency and diverse motion, and Hunyuan Video~\citep{kong2024hunyuanvideo}, which employs a causal 3D VAE~\citep{kingma2013auto} for spatio-temporal token compression in latent space combined with a large language model for text conditioning. Wan 2.1~\citep{wan2025} further demonstrates the benefits of massive pretraining for high-resolution video generation. CogVideoX~\citep{hong2022cogvideo,yang2024cogvideox} introduces an expert transformer with adaptive LayerNorm to enhance cross-modal fusion, supported by a 3D VAE, progressive training, and multi-resolution frame packing, thereby achieving strong text alignment and motion coherence. Open-Sora~\citep{lin2024open,peng2025open} and Open-Sora-Plan~\citep{lin2024open} extend these advancements in the open-source community, delivering high-quality video generation and significantly accelerating progress in efficiency and realism. Collectively, these works illustrate how scaling strategies and architectural innovations have transformed video diffusion from modest UNet adaptations into transformer-driven models capable of generating controllable and high-quality videos.


\textbf{Long Video Generation}
Due to the substantial training and inference cost of DiT-based architectures, most state-of-the-art models remain limited to generating videos of 5–10 seconds. To overcome this constraint, a number of techniques have been introduced to extend generation to longer durations~\citep{jiang2025lovic,kodaira2025streamdit,fang2025inflvg,lin2025autoregressive}. RIFLEx~\citep{zhao2025riflex} is a training-free approach that revisits positional encoding, effectively doubling the generation length by avoiding encodings that induce repetitive motion, and surpassing prior methods by a large margin~\citep{chen2023extending,peng2023yarn,zhuo2024lumina}. Another promising direction is autoregressive video generation. Nova~\citep{deng2024autoregressive} reformulates video synthesis as a non-quantized autoregressive problem, jointly modeling temporal frame-by-frame prediction and spatial set-by-set prediction, which enables flexible in-context learning. Pyramid-Flow~\citep{jin2024pyramidal} interprets denoising as a hierarchical process across multi-stage pyramids, linking flows across resolutions and time to support end-to-end autoregressive video generation with a single diffusion transformer. SkyReels-V2 adopts diffusion forcing~\citep{chen2024diffusion} to support potentially infinite rollouts, while MAGI-1~\citep{teng2025magi} trains a model to progressively denoise per-chunk noise that increases over time, autoregressively predicting fixed-length segments of consecutive frames. CausVid~\citep{yin2025causvid} employs block causal attention and a KV cache to autoregressively extend sequences, and Self-Forcing~\citep{huang2025self} further aligns training with inference by incorporating the KV cache directly during training, producing high-quality short videos.

\textbf{Reinforcement Learning}
Reinforcement learning has become a central component in the post-training of large language models~\citep{ouyang2022training,kumar2024training,hou2025advancing,xie2025teaching}. With the rise of image generation, it has also proven effective for improving generative models, with early efforts introducing reward models tailored to images, such as ImageReward~\citep{xu2023imagereward}, Pick-a-Pic~\citep{kirstain2023pick}, and HPS V2~\citep{wu2023human}. These concepts have since been extended to video through reward functions like VideoReward~\citep{liu2025improving} and VisionReward~\citep{xu2024visionrewardfinegrainedmultidimensionalhuman}, which assess temporal coherence and motion quality. Building on these reward signals, optimization techniques first developed for language models, including Direct Preference Optimization (DPO)\citep{rafailov2023direct} and Group Relative Policy Optimization (GRPO)\citep{shao2024deepseekmath}, have been adapted to diffusion-based generation. Notably, Diffusion-DPO~\citep{wallace2024diffusion} applies preference-based training directly to diffusion models, while Flow-GRPO~\citep{liu2025flow,xue2025dancegrpo} leverages GRPO to fine-tune video diffusion models, resulting in improvements in both visual fidelity and motion consistency.